% GENERATED FILE. Edit canonical lessons, then run npm run generate:books.
% canonical-source-hash: fnv1a64:6f17b6ad8768c6df
% canonical-lessons: BN-C531-khosra, BN-C531-talika, BN-C531-krom, BN-C531-ruprekha, BN-C531-bad-deoya, BN-R531-review-524-531

\chapter{Drafts and Lists}
\label{ch:bn-drafts-and-lists}
\hlchaptermodality{531}

\begin{chapteropening}
\textbf{By the end of this chapter:} \emph{I can talk about drafts, lists and working out a plan.}

Complete the last lesson of chapter 531: I can talk about drafts, lists and working out a plan.
\end{chapteropening}

\section[\texorpdfstring{\bn{খসড়া} (khôsṛā)}{খসড়া (khôsṛā)}]{\bn{খসড়া} (khôsṛā) — a draft (a first version)}

\label{lesson:BN-C531-khosra}



\begin{quote}
Before the new one: say the Bengali for to stall, then the Bengali for to settle on.
\end{quote}

\subsection*{You'll want to know: \bn{খসড়া}}
\textbf{\bn{খসড়া}} — \emph{khôsṛā} — \textquotedblleft{}a draft (a first version)\textquotedblright{}.

\begin{quote}
\bn{আমি} \bn{পরিকল্পনার} \bn{একটা} \bn{খসড়া} \bn{বানালাম।}
\end{quote}

\begin{quote}
\emph{Āmi pôrikôlpônār ekṭā khôsṛā bānālām.}
\end{quote}

\begin{quote}
I made a draft of the plan.
\end{quote}

\begin{cousinweb}
\textbf{\bn{খসড়া}} (\emph{khôsṛā}) is usually traced to an Arabic word \emph{khasra}, a rough record, which came through Persian.
\end{cousinweb}

\begin{grammarlens}[title={What you've built}]
One more word for saying what happens if.
\end{grammarlens}

\subsection*{Your turn}
\begin{itemize}
\raggedright
  \item \emph{Say it:} \emph{khôsṛā}
  \item \emph{Say it:} \emph{khôsṛā}, once more
  \item \emph{Say it:} \emph{Āmi pôrikôlpônār ekṭā khôsṛā bānālām.} — I made a draft of the plan.
  \item \emph{Recall:} say the Bengali for to stall, then the Bengali for to settle on, then say \emph{khôsṛā} again
\end{itemize}

\begin{tcolorbox}[breakable,colback=teal!4,colframe=teal!35!black,title={Before you move on}]
What does \bn{খসড়া} mean? (\textquotedblleft{}A draft (a first version)\textquotedblright{}.) Say it once more.
\end{tcolorbox}
\section[\texorpdfstring{\bn{তালিকা} (tālikā)}{তালিকা (tālikā)}]{\bn{তালিকা} (tālikā) — a list}

\label{lesson:BN-C531-talika}



\begin{quote}
Before the new one: say the Bengali for to settle on, then the Bengali for a draft.
\end{quote}

\subsection*{You'll want to know: \bn{তালিকা}}
\textbf{\bn{তালিকা}} — \emph{tālikā} — \textquotedblleft{}a list\textquotedblright{}.

\begin{quote}
\bn{বাজারে} \bn{যাওয়ার} \bn{আগে} \bn{আমি} \bn{একটা} \bn{তালিকা} \bn{লিখলাম।}
\end{quote}

\begin{quote}
\emph{Bājāre jāoyār āge āmi ekṭā tālikā likhlām.}
\end{quote}

\begin{quote}
Before going to the market I wrote a list.
\end{quote}

\begin{cousinweb}
\textbf{\bn{তালিকা}} (\emph{tālikā}) is usually traced to Arabic \emph{taʿlīqa}, \textquotedblleft{}a note, an annotation\textquotedblright{}, through Persian.
\end{cousinweb}

\begin{grammarlens}[title={What you've built}]
One more word for saying what happens if.
\end{grammarlens}

\subsection*{Your turn}
\begin{itemize}
\raggedright
  \item \emph{Say it:} \emph{tālikā}
  \item \emph{Say it:} \emph{tālikā}, once more
  \item \emph{Say it:} \emph{Bājāre jāoyār āge āmi ekṭā tālikā likhlām.} — Before going to the market I wrote a list.
  \item \emph{Recall:} say the Bengali for to settle on, then the Bengali for a draft, then say \emph{tālikā} again
\end{itemize}

\begin{tcolorbox}[breakable,colback=teal!4,colframe=teal!35!black,title={Before you move on}]
What does \bn{তালিকা} mean? (\textquotedblleft{}A list\textquotedblright{}.) Say it once more.
\end{tcolorbox}
\section[\texorpdfstring{\bn{ক্রম} (krôm)}{ক্রম (krôm)}]{\bn{ক্রম} (krôm) — an order; a sequence}

\label{lesson:BN-C531-krom}



\begin{quote}
Before the new one: say the Bengali for a draft, then the Bengali for a list.
\end{quote}

\subsection*{You'll want to know: \bn{ক্রম}}
\textbf{\bn{ক্রম}} — \emph{krôm} — \textquotedblleft{}an order; a sequence\textquotedblright{}.

\begin{quote}
\bn{আমি} \bn{ফাইলগুলো} \bn{তারিখের} \bn{ক্রমে} \bn{সাজালাম।}
\end{quote}

\begin{quote}
\emph{Āmi phāilgulo tārikher krôme sājālām.}
\end{quote}

\begin{quote}
I arranged the files in order of date.
\end{quote}

\begin{cousinweb}
\textbf{\bn{ক্রম}} (\emph{krôm}) is a tatsama, Sanskrit \emph{krama}, \textquotedblleft{}a step\textquotedblright{}, and from that \textquotedblleft{}regular order, succession\textquotedblright{} (Monier-Williams).
\end{cousinweb}

\begin{grammarlens}[title={What you've built}]
One more word for saying what happens if.
\end{grammarlens}

\subsection*{Your turn}
\begin{itemize}
\raggedright
  \item \emph{Say it:} \emph{krôm}
  \item \emph{Say it:} \emph{krôm}, once more
  \item \emph{Say it:} \emph{Āmi phāilgulo tārikher krôme sājālām.} — I arranged the files in order of date.
  \item \emph{Recall:} say the Bengali for a draft, then the Bengali for a list, then say \emph{krôm} again
\end{itemize}

\begin{tcolorbox}[breakable,colback=teal!4,colframe=teal!35!black,title={Before you move on}]
What does \bn{ক্রম} mean? (\textquotedblleft{}An order; a sequence\textquotedblright{}.) Say it once more.
\end{tcolorbox}
\section[\texorpdfstring{\bn{রূপরেখা} (rūprekhā)}{রূপরেখা (rūprekhā)}]{\bn{রূপরেখা} (rūprekhā) — an outline}

\label{lesson:BN-C531-ruprekha}



\begin{quote}
Before the new one: say the Bengali for a list, then the Bengali for a sequence.
\end{quote}

\subsection*{You'll want to know: \bn{রূপরেখা}}
\textbf{\bn{রূপরেখা}} — \emph{rūprekhā} — \textquotedblleft{}an outline\textquotedblright{}.

\begin{quote}
\bn{ম্যানেজার} \bn{নতুন} \bn{কাজের} \bn{একটা} \bn{রূপরেখা} \bn{দিলেন।}
\end{quote}

\begin{quote}
\emph{Mænejār nôtun kājer ekṭā rūprekhā dilen.}
\end{quote}

\begin{quote}
The manager gave an outline of the new work.
\end{quote}

\begin{cousinweb}
\textbf{\bn{রূপরেখা}} (\emph{rūprekhā}) joins Sanskrit \emph{rūpa}, \textquotedblleft{}form, shape\textquotedblright{}, and \emph{rekhā}, \textquotedblleft{}a line, a stroke\textquotedblright{} (Monier-Williams): the lines of a shape before it is filled in.
\end{cousinweb}

\begin{grammarlens}[title={What you've built}]
One more word for saying what happens if.
\end{grammarlens}

\subsection*{Your turn}
\begin{itemize}
\raggedright
  \item \emph{Say it:} \emph{rūprekhā}
  \item \emph{Say it:} \emph{rūprekhā}, once more
  \item \emph{Say it:} \emph{Mænejār nôtun kājer ekṭā rūprekhā dilen.} — The manager gave an outline of the new work.
  \item \emph{Recall:} say the Bengali for a list, then the Bengali for a sequence, then say \emph{rūprekhā} again
\end{itemize}

\begin{tcolorbox}[breakable,colback=teal!4,colframe=teal!35!black,title={Before you move on}]
What does \bn{রূপরেখা} mean? (\textquotedblleft{}An outline\textquotedblright{}.) Say it once more.
\end{tcolorbox}
\section[\texorpdfstring{\bn{বাদ} \bn{দেওয়া} (bād deoyā)}{বাদ দেওয়া (bād deoyā)}]{\bn{বাদ} \bn{দেওয়া} (bād deoyā) — to leave out; to drop (from a list)}

\label{lesson:BN-C531-bad-deoya}



\begin{quote}
Before the new one: say the Bengali for a sequence, then the Bengali for an outline.
\end{quote}

\subsection*{You'll want to know: \bn{বাদ} \bn{দেওয়া}}
\textbf{\bn{বাদ} \bn{দেওয়া}} — \emph{bād deoyā} — \textquotedblleft{}to leave out; to drop (from a list)\textquotedblright{}.

\begin{quote}
\bn{এই} \bn{নামটা} \bn{বাদ} \bn{দিলে} \bn{তালিকাটা} \bn{ছোট} \bn{হবে।}
\end{quote}

\begin{quote}
\emph{Ei nāmṭā bād dile tālikāṭā chhoṭo hôbe.}
\end{quote}

\begin{quote}
If you drop this name, the list will be shorter.
\end{quote}

Past (he/she …): \bn{তিনি} \bn{পুরনো} \bn{খরচটা} \bn{বাদ} \bn{দিলেন}, \emph{tini purôno khôrôchṭā bād dilen}.

\begin{cousinweb}
\textbf{\bn{বাদ} \bn{দেওয়া}} (\emph{bād deoyā}) joins \bn{বাদ}, \textquotedblleft{}omission, leaving aside\textquotedblright{}, usually traced to Arabic \emph{baʿd}, \textquotedblleft{}after\textquotedblright{}, through Persian, to \bn{দেওয়া}, \textquotedblleft{}to give\textquotedblright{}; the sense \textquotedblleft{}leaving out\textquotedblright{} is not the Arabic one.
\end{cousinweb}

\begin{grammarlens}[title={What you've built}]
One more word for saying what happens if.
\end{grammarlens}

\subsection*{Your turn}
\begin{itemize}
\raggedright
  \item \emph{Say it:} \emph{bād deoyā}
  \item \emph{Say it:} \emph{bād deoyā}, once more
  \item \emph{Say it:} \emph{Ei nāmṭā bād dile tālikāṭā chhoṭo hôbe.} — If you drop this name, the list will be shorter.
  \item \emph{Recall:} say the Bengali for a sequence, then the Bengali for an outline, then say \emph{bād deoyā} again
\end{itemize}

\begin{tcolorbox}[breakable,colback=teal!4,colframe=teal!35!black,title={Before you move on}]
What does \bn{বাদ} \bn{দেওয়া} mean? (\textquotedblleft{}To leave out; to drop (from a list)\textquotedblright{}.) Say it once more.
\end{tcolorbox}
\section[(dialogue)]{Looking back: from Paint and Carving to Drafts and Lists}

\label{lesson:BN-R531-review-524-531}



\begin{quote}
Say the Bengali for an easel and watercolour.
\end{quote}

\subsection*{Your turn}
Say these aloud:

\begin{itemize}
\raggedright
  \item an easel, watercolour, a pattern, to carve, to dye
  \item a tune, rhythm, musical notation, to resound, to whistle
  \item a sitar, a drum, a harmonium, a violin, to strum
  \item a stage, to act on stage, an actor, a rehearsal, applause
  \item financial means, a way out, capacity, to suffice, to scrape together
  \item pleasing to oneself, a dilemma, personal taste, to lean forward, a whim
  \item adjourned, fixed in advance, in advance, to stall, to settle on
  \item a draft, a list, a sequence, an outline, to leave out
\end{itemize}

\begin{tcolorbox}[breakable,colback=teal!4,colframe=teal!35!black,title={Before you move on}]
An easel? (\textbf{\bn{ইজেল}}, \emph{ijel}.) Watercolour? (\textbf{\bn{জলরং}}, \emph{jôlrông}.) A pattern? (\textbf{\bn{নকশা}}, \emph{nôkshā}.) To carve? (\textbf{\bn{খোদাই} \bn{করা}}, \emph{khodāi kôrā}.) To dye? (\textbf{\bn{রাঙানো}}, \emph{rāngāno}.) A tune? (\textbf{\bn{সুর}}, \emph{sur}.) Rhythm? (\textbf{\bn{তাল}}, \emph{tāl}.) Musical notation? (\textbf{\bn{স্বরলিপি}}, \emph{svôrlipi}.) To resound? (\textbf{\bn{গমগম} \bn{করা}}, \emph{gômgôm kôrā}.) To whistle? (\textbf{\bn{শিস} \bn{দেওয়া}}, \emph{shis deoyā}.) A sitar? (\textbf{\bn{সেতার}}, \emph{setār}.) A drum? (\textbf{\bn{ঢোল}}, \emph{ḍhol}.) A harmonium? (\textbf{\bn{হারমোনিয়াম}}, \emph{hārmoniyām}.) A violin? (\textbf{\bn{বেহালা}}, \emph{behālā}.) To strum? (\textbf{\bn{টুংটাং} \bn{করা}}, \emph{ṭungṭāng kôrā}.) A stage? (\textbf{\bn{মঞ্চ}}, \emph{mônchô}.) To act on stage? (\textbf{\bn{অভিনয়} \bn{করা}}, \emph{ôbhinôy kôrā}.) An actor? (\textbf{\bn{অভিনেতা}}, \emph{ôbhinetā}.) A rehearsal? (\textbf{\bn{রিহার্সাল}}, \emph{rihārsāl}.) Applause? (\textbf{\bn{হাততালি}}, \emph{hāttāli}.) Financial means? (\textbf{\bn{সঙ্গতি}}, \emph{sônggôti}.) A way out? (\textbf{\bn{ফিকির}}, \emph{phikir}.) Capacity? (\textbf{\bn{সামর্থ্য}}, \emph{sāmôrthyô}.) To suffice? (\textbf{\bn{কুলোনো}}, \emph{kulono}.) To scrape together? (\textbf{\bn{জোটানো}}, \emph{joṭāno}.) Pleasing to oneself? (\textbf{\bn{মনঃপূত}}, \emph{mônôppūtô}.) A dilemma? (\textbf{\bn{দোটানা}}, \emph{doṭānā}.) Personal taste? (\textbf{\bn{রুচি}}, \emph{ruchi}.) To lean forward? (\textbf{\bn{ঝোঁকা}}, \emph{jhõkā}.) A whim? (\textbf{\bn{মর্জি}}, \emph{môrji}.) Adjourned? (\textbf{\bn{মুলতুবি}}, \emph{multubi}.) Fixed in advance? (\textbf{\bn{নির্দিষ্ট}}, \emph{nirdishṭô}.) In advance? (\textbf{\bn{আগাম}}, \emph{āgām}.) To stall? (\textbf{\bn{টালবাহানা} \bn{করা}}, \emph{ṭālbāhānā kôrā}.) To settle on? (\textbf{\bn{সাব্যস্ত} \bn{করা}}, \emph{sābyôstô kôrā}.) A draft? (\textbf{\bn{খসড়া}}, \emph{khôsṛā}.) A list? (\textbf{\bn{তালিকা}}, \emph{tālikā}.) A sequence? (\textbf{\bn{ক্রম}}, \emph{krôm}.) An outline? (\textbf{\bn{রূপরেখা}}, \emph{rūprekhā}.) To leave out? (\textbf{\bn{বাদ} \bn{দেওয়া}}, \emph{bād deoyā}.)
\end{tcolorbox}
